\begin{pastexam}[label={px:ncku-106-prog-03}]{成大資訊聯招 106 學年度 第 3 題}
（問答）

(10\%) What are the minimum and maximum numbers of elements in a heap of height $h$?
\end{pastexam}

\begin{pastsoln}
題目沒有定義高度。 這題的說法來自英文教科書（CLRS），那裡的高度數的是\textbf{邊數}：只有根的堆積高度是
$0$。 資結書第~7.1~節的高度數的是\textbf{節點數}（層數），只有根的高度是 $1$。 下面兩種都寫，作答時先寫
一句「這裡的高度是邊數」。

答案：
\begin{center}
\begin{tabular}{lcc}
\toprule
高度的定義 & 最少 & 最多 \\
\midrule
數邊數（根的深度 $0$，最深的節點深度 $h$） & $2^h$ & $2^{h+1} - 1$ \\
數節點數（共 $h$ 層） & $2^{h-1}$ & $2^h - 1$ \\
\bottomrule
\end{tabular}
\end{center}

\textbf{推導（高度數邊數）。} 二元堆積是完全二元樹（資結書第~8.1~節）：除了最深的一層，每一層都排滿，
最深的一層從左邊排起。 深度 $d$ 的一層排滿時有 $2^d$ 個節點（每往下一層，個數加倍）。 高度是 $h$，表示
最深的一層是深度 $h$，而且至少有一個節點。
\begin{itemize}
\item 深度 $0$ 到 $h - 1$ 的 $h$ 層都排滿，共 $\sum_{d=0}^{h-1} 2^d = 2^h - 1$ 個節點。
\item 深度 $h$ 那一層有 $1$ 到 $2^h$ 個節點。
\end{itemize}
所以元素個數 $n$ 滿足
\[
(2^h - 1) + 1 = 2^h \;\le\; n \;\le\; (2^h - 1) + 2^h = 2^{h+1} - 1 .
\]
兩端都達得到：最深一層只有最左邊一個節點時是 $2^h$，最深一層排滿時是 $2^{h+1} - 1$。 這也說明 $n$ 個元素
的堆積高度是 $\lfloor \lg n \rfloor$：$2^h \le n < 2^{h+1}$。

\textbf{高度數節點數時}，高度 $h$ 的堆積有 $h$ 層，最深的節點深度 $h - 1$，把上面的 $h$ 換成 $h - 1$：
最少 $2^{h-1}$、最多 $2^h - 1$。

例：邊數高度 $h = 2$ 時，最少 $4$ 個（前兩層 $3$ 個，第三層 $1$ 個），最多 $7$ 個（三層排滿）。
\end{pastsoln}
